\section{Residue characteristic zero}\label{sec:charzero}

We deduce the residue-characteristic-zero case from
Theorem~\ref{thm:positive-characteristic} by encoding a hypothetical
strict counterexample in finitely many polynomial equations. We follow
the reduction-and-approximation strategy of \cite[Section~5]{Ma},
beginning with the finite free reduction of \cite[Lemma~5.1]{Ma}.
For the subsequent transfer, it is enough to preserve a lower bound for
the source multiplicity and an upper bound for the target multiplicity. This
allows us to use a presentation matrix and finitely many initial
relations, without keeping track of the homology of a resolution.

\subsection{A finite free map with a common coefficient field}

\begin{lemma}\label{lem:charzero-finite-free}
Let $(R,\mathfrak m)\to(S,\mathfrak n)$ be a flat local map between
complete Noetherian local rings of equal dimension and residue
characteristic zero.  There is a factorization
\[
 R\longrightarrow R'\longrightarrow S
\]
such that $R'$ is complete local, $e(R')=e(R)$, and $R'\to S$ is finite
free.  Moreover, $R'$ and $S$ have the same coefficient field and the
same residue field.
\end{lemma}

\begin{proof}
Choose a coefficient field $k\subset R$, and identify it with its image
in $S$.  Put $L=S/\mathfrak n$.  A coefficient field of $S$ containing
this copy of $k$ can be chosen as follows.  Lift a transcendence basis
of $L/k$ to $S$.  These lifts generate a copy of the corresponding
rational function field, because a polynomial with nonzero residue is
a unit whenever it occurs as a denominator.  The remaining algebraic
extension of residue fields is separable.  Hensel's lemma lifts each
simple algebraic root, and a maximal compatible choice, obtained by
Zorn's lemma, lifts all of $L$.  We henceforth identify this coefficient
field with $L$.

Choose a presentation $R=k[[X_1,\ldots,X_a]]/I$ and set
\[
 R'=L[[X_1,\ldots,X_a]]/I L[[X_1,\ldots,X_a]].
\]
The compatible maps from $R$ and $L$ to $S$ give a local map $R'\to S$.
Extension of coefficient fields on power series rings is flat; for
example, the local criterion of flatness applies to the regular
sequence of variables and the field extension on their quotients.
Thus $R\to R'$ is flat, its maximal ideal is $\mathfrak m'=\mathfrak mR'$,
and its residue field is $L$.  Filtering by powers of the maximal
ideal gives
\[
 \gr_{\mathfrak m'}R'
   \cong (\gr_{\mathfrak m}R)\otimes_k L.
\]
In particular, $\dim R'=\dim R$ and $e(R')=e(R)$.

The flat dimension formula for the original map shows that
$S/\mathfrak mS$ is Artinian.  Since
$\mathfrak m'S=\mathfrak mS$, this quotient is a finite-dimensional
$L$-vector space.  Lift an $L$-basis to $s_1,\ldots,s_q\in S$.
Successively expressing a remainder modulo $\mathfrak m'S$, then
modulo $(\mathfrak m'S)^2$, and so on, expresses every element of $S$
as $\sum_i a_i s_i$ with $a_i\in R'$.  Indeed, the coefficient
corrections at step $j$ lie in $(\mathfrak m')^j$, so they converge in
$R'$, and the remainders tend to zero in $S$: the
$\mathfrak m'S$-adic and $\mathfrak n$-adic topologies agree.  Hence
$S$ is finite over $R'$.

Let $P_\bullet$ be a free resolution of $k$ over $R$.  Flatness of
$R'$ over $R$ makes $P_\bullet\otimes_R R'$ a free resolution of $L$
over $R'$.  Therefore
\[
 \Tor^{R'}_1(L,S)
   = H_1(P_\bullet\otimes_R S)
   = \Tor^R_1(k,S)=0.
\]
A finite module over a Noetherian local ring with this vanishing is
free.  Thus $R'\to S$ is finite free, necessarily of positive rank.
\end{proof}

\subsection{Encoding the multiplicity bounds}

\begin{lemma}\label{lem:charzero-certificate}
Let $(R_0,\mathfrak m_0)\to(S_0,\mathfrak n_0)$ be a finite free
local map of complete local rings of dimension $d>0$, with a common
infinite coefficient field $L$ and residue field $L$.  Put
$r=e(R_0)$.  There is a finite system
\[
 F(t_1,\ldots,t_d,Y_1,\ldots,Y_N)=0
\]
of polynomial equations over $\ZZ$ with a solution in
$L[[t_1,\ldots,t_d]]$, having the following property.  For every field
$K$ and every solution in $A=K[[t_1,\ldots,t_d]]$, the solution
determines a finite free local map
\[
 (R,\mathfrak a)\longrightarrow(S,\mathfrak b)
\]
of nonzero complete local rings of dimension $d$ and residue field
$K$, for which
\[
 e(R)\ge r,\qquad e(S)\le e(S_0).
\]
\end{lemma}

\begin{proof}
Choose a minimal reduction $t_1,\ldots,t_d$ of $\mathfrak m_0$.
The induced map
\[
 A_0=L[[t_1,\ldots,t_d]]\longrightarrow R_0
\]
is finite, by completeness and the parameter property.  It is
injective: a nonzero kernel in the regular local domain $A_0$ would
force the dimension of its finite algebra $R_0$ to be less than $d$.
Since the residue fields agree, the finite-module multiplicity formula
and the reduction property give
\begin{equation}\label{eq:charzero-original-rank}
 r=e(R_0)=e((t),R_0)=\rank_{A_0}R_0.
\end{equation}
Here no domain or Cohen--Macaulay hypothesis on $R_0$ is needed.
For completeness, a finite module of rank $r$ over $A_0$ contains a
copy of $A_0^r$ with torsion cokernel; the cokernel has dimension less
than $d$, so additivity of multiplicity gives the last equality.

We describe the finite system.  Throughout the construction, an
equality in a presented module means a coordinate equality with an
additional vector witnessing membership in its relation submodule.
Thus, for a matrix $U$, the assertion $w\in\operatorname{im}U$ is
encoded by $w=Uc$ with a new vector of unknowns $c$.  The number of
coordinates and assertions below is always fixed and finite.

The system records the algebra structures, freeness, and locality.
The numerical conditions use the source's generic rank for
a lower multiplicity bound and the dimensions of the target's associated
graded pieces for an upper bound.

\smallskip
\noindent\emph{The source algebra and its generic rank.}
Fix a presentation
\[
 A_0^v\xrightarrow{U_0}A_0^u\longrightarrow R_0\longrightarrow0.
\]
The entries of a matrix $U$ of this same size are unknowns.  On
$A^u$, specify a bilinear multiplication by vectors $c_{ij}\in A^u$
giving the products of standard basis vectors, and specify a unit
vector $\epsilon\in A^u$.  Impose the following identities modulo
$\operatorname{im}U$: commutativity and associativity on basis
vectors, the two unit identities on basis vectors, and the assertion
that a relation column multiplied by any basis vector is again in
$\operatorname{im}U$.  These identities give a unital commutative
$A$-algebra structure on
\[
 R=\operatorname{coker}U.
\]
They are polynomial equations in the entries of $U$, the
multiplication table, $\epsilon$, and the membership witnesses.

Impose also the vanishing of all $(u-r+1)$-minors of $U$.  If there
are no minors of that size, this condition is empty.  Over the
fraction field of $A$, it gives
\begin{equation}\label{eq:charzero-rank-bound}
 \rank_A R=u-\rank U\ge r>0.
\end{equation}
These equations are satisfied by the original presentation in view of
\eqref{eq:charzero-original-rank}. A solution may give $U$ smaller rank,
but this only increases the rank of its cokernel. To turn that rank
bound into a multiplicity bound, we will retain the source reduction
$(t)$ as well.

\smallskip
\noindent\emph{The target as a free source-module.}
Let $n>0$ be the original free rank of $S_0$ over $R_0$.
Define the underlying module of the new target to be $S=R^n$.
Give it an $R$-bilinear multiplication by specifying
\[
 \eta_i\eta_j=\sum_{h=1}^n d_{ijh}\eta_h,
 \qquad d_{ijh}\in R,
\]
and a unit $\delta\in R^n$.  Represent these coefficients and this
unit by vectors over $A$.  Impose commutativity, associativity, and
the unit identities on the $R$-basis $\eta_1,\ldots,\eta_n$.
The underlying $A$-module is presented by the block diagonal matrix
with $n$ copies of $U$, so these are again finitely many polynomial
equations, using the source multiplication table and relation
witnesses.  The structural map is $a\mapsto a\delta$.
Consequently, in every solution, $S$ is a unital commutative
$R$-algebra which is free of rank $n$ as an $R$-module.

\smallskip
\noindent\emph{The local ideals and the source reduction.}
Choose finite generating lists for $\mathfrak m_0$ and
$\mathfrak n_0$ that include the respective images of $t_1,\ldots,t_d$.
Represent corresponding lists in $R$ and $S$ by vectors of
unknowns, with their first $d$ entries defined to be
$t_i\epsilon$ and $t_i\delta$.  Denote the generated ideals by
$\mathfrak a$ and $\mathfrak b$.  Fix exponents $N_R,N_S$ and $j$
for which the original rings satisfy the following conditions, and
impose them on the new data:
\begin{equation}\label{eq:charzero-ideal-conditions}
 \mathfrak a^{N_R}\subseteq(t)R,\qquad
 \mathfrak b^{N_S}\subseteq(t)S,\qquad
 \mathfrak a^{j+1}\subseteq(t)\mathfrak a^j.
\end{equation}
Such exponents exist because $(t)$ is primary for both original
maximal ideals and reduces $\mathfrak m_0$.
For each standard $A$-module generator $e_i$ of $R$, impose an
expression
\[
 e_i=c_i\epsilon+\sum_\nu a_\nu w_{i\nu},
 \qquad c_i\in A,\quad w_{i\nu}\in R,
\]
where $a_\nu$ is the source ideal list.  Impose the analogous
expressions for the $un$ standard $A$-module generators of $S$,
using its unit and its ideal list.  Thus $R/\mathfrak a$ and
$S/\mathfrak b$ are spanned by their units as $A$-modules.
All these conditions reduce to identities on finite lists of
generators and monomials.  Expanding products with the multiplication
tables, and representing each coefficient by a vector, encodes them
as polynomial equations.

We first check their consequences, before imposing the target
Hilbert-function conditions.  By \eqref{eq:charzero-rank-bound}, $R$
is a nonzero finite $A$-algebra, and $S=R^n$ is also nonzero and
finite over $A$.  In a finite algebra over the local ring $A$, every
maximal ideal contains the extended ideal $(t)$.  The first two
conditions in \eqref{eq:charzero-ideal-conditions} therefore show
that every maximal ideal of $R$ contains $\mathfrak a$, and every
maximal ideal of $S$ contains $\mathfrak b$.  In particular, both
ideals are proper.  Since $R/\mathfrak a$ is spanned by the unit and
annihilated by $(t)$, it is a nonzero quotient of $K$, hence is $K$.
Thus $\mathfrak a$ is the unique maximal ideal of $R$.  The same
argument gives $S/\mathfrak b=K$ and uniqueness of $\mathfrak b$.

The map $R\to S$ is local.  Indeed, an element of $\mathfrak a$
is nilpotent modulo $(t)R$, so its image belongs to
$\sqrt{(t)S}=\mathfrak b$; contraction of $\mathfrak b$ is then
$\mathfrak a$.  Positive generic rank implies that $R$ is faithful
as an $A$-module, so $A\to R$ is injective.  The map $R\to S$ is
also injective: if $a\delta=0$, multiplication by $a$ on $S$ is
zero, whereas $S=R^n$ is faithful.  Finiteness now gives
$\dim R=\dim S=d$.  Both rings are complete, since they are finite
over complete $A$ and their maximal-ideal topologies agree with their
$(t)$-adic topologies by \eqref{eq:charzero-ideal-conditions}.

Because $(t)\subseteq\mathfrak a$, the last inclusion in
\eqref{eq:charzero-ideal-conditions} is equality.  It makes $(t)$
a reduction of $\mathfrak a$, whence
\begin{equation}\label{eq:charzero-source-bound}
 e(R)=e((t),R)=\rank_A R\ge r.
\end{equation}
The rank equality has the same proof as
\eqref{eq:charzero-original-rank}. Thus the reduction relation converts
the lower rank bound into the required source multiplicity bound.
Neither a nonvanishing minor nor an additional condition on the length
of the residue quotient is required.

\smallskip
\noindent\emph{An upper bound for the target Hilbert function.}
Write $y_1,\ldots,y_q$ for the chosen original target ideal list.
There is a graded surjection
\[
 L[Z_1,\ldots,Z_q]\longrightarrow\gr_{\mathfrak n_0}S_0,
 \qquad Z_i\longmapsto y_i+\mathfrak n_0^2.
\]
Let $I$ be its homogeneous kernel.  Fix a degree-compatible monomial
order.  Choose finitely many monic homogeneous polynomials
$f_1,\ldots,f_b\in I$ whose leading monomials generate the initial
ideal $J=\operatorname{in}(I)$.  Write
\[
 f_\ell=Z^{\alpha_\ell}
       +\sum_{\beta\in B_\ell}c_{\ell\beta}Z^\beta,
 \qquad |\beta|=|\alpha_\ell|=a_\ell,
 \quad Z^\beta<Z^{\alpha_\ell}.
\]
Here $B_\ell$ is a fixed finite set.  Replace the coefficients
$c_{\ell\beta}$ by unknowns $C_{\ell\beta}\in A$ and require, on
the new target ideal list $b_1,\ldots,b_q$, that
\begin{equation}\label{eq:charzero-leading-relations}
 b^{\alpha_\ell}
  +\sum_{\beta\in B_\ell}C_{\ell\beta}b^\beta
       \in\mathfrak b^{a_\ell+1}.
\end{equation}
Membership is expressed using the finitely many degree
$a_\ell+1$ monomials in the ideal list and coefficient vectors in
$S$.  These are again polynomial equations, and the original data
provide a solution.

For any transferred solution, reduce the coefficients
$C_{\ell\beta}$ modulo $(t)$.  Equation
\eqref{eq:charzero-leading-relations} gives a homogeneous relation in
the kernel of $K[Z]\to\gr_{\mathfrak b}S$ with the same leading
monomial $Z^{\alpha_\ell}$.  The leading coefficient is still $1$.
The higher $(t)$-adic terms of the other coefficients do not contribute:
their products with a degree-$a_\ell$ monomial lie in
$\mathfrak b^{a_\ell+1}$.  It follows that the new initial ideal
contains the monomial ideal $J$, interpreted over $K$. The transferred
target may have additional initial relations; this containment is the
direction needed for an upper bound. Standard monomials outside $J$
thus span each new graded piece, whereas they
form a basis of the corresponding original graded piece.  Consequently
\[
 \dim_K\mathfrak b^a/\mathfrak b^{a+1}
 \le \dim_L\mathfrak n_0^a/\mathfrak n_0^{a+1}
 \quad(a\ge0).
\]
Summation and the equality of dimensions give
\begin{equation}\label{eq:charzero-target-bound}
 e(S)\le e(S_0).
\end{equation}

Finally, every coefficient in the foregoing presentations, tables,
lists, and membership witnesses is an unknown; only the sizes,
exponents, monomials, and the distinguished coefficients $1$ have
been fixed.  Accordingly the entire system has coefficients in
$\ZZ[t_1,\ldots,t_d]$.  It is finite, the original data solve it in
$A_0$, and \eqref{eq:charzero-source-bound} and
\eqref{eq:charzero-target-bound} prove the assertion.
\end{proof}

\subsection{Specialization to positive characteristic}

The preceding encoding reduces the problem to transferring existence
of a solution of a finite polynomial system. Artin approximation supplies
an algebraic solution whose coefficients can be specialized.

\begin{lemma}\label{lem:charzero-specialization}
Let $L$ be a field of characteristic zero.  A finite polynomial system
over $\ZZ[t_1,\ldots,t_d]$ which has a solution in $L[[t_1,\ldots,t_d]]$
has a solution in $K[[t_1,\ldots,t_d]]$ for some algebraically closed
field $K$ of positive characteristic.
\end{lemma}

\begin{proof}
Put $B=(L[t_1,\ldots,t_d]_{(t)})^h$, whose completion is $L[[t]]$.
Artin approximation \cite[Theorem~1.10]{Artin} gives a solution in
$B$.  The henselization is the filtered direct limit of pointed
etale neighborhoods of the origin.  Thus the finitely many solution
entries and identities descend to one such neighborhood with an
$L$-rational point over $t=0$; this is also the explicit conclusion
of \cite[Corollary~2.1]{Artin}.

Shrink about that point.  An etale presentation then has the form
\[
 E=\bigl(L[t,z_1,\ldots,z_m]/(G_1,\ldots,G_m)\bigr)_{g\Delta},
 \qquad
 \Delta=\det\left(\frac{\partial G_i}{\partial z_j}\right),
\]
where the distinguished point is $(t,z)=(0,c)$, with
$c\in L^m$, $G_i(0,c)=0$, and $g(0,c)\Delta(0,c)\ne0$.
The solution entries are elements of $E$.  Choose polynomial
representatives with denominators powers of $g\Delta$.  For each
of the finitely many solution identities, clear denominators and
choose a polynomial expression in the ideal $(G_1,\ldots,G_m)$
which witnesses that identity.  One may multiply by a further power
of $g\Delta$ before doing so.

Let $W\subset L$ be the $\ZZ$-algebra generated by the following
finite set: the coefficients of $G_i$ and $g$, the coordinates of
$c$, the coefficients of the representatives and ideal-membership
witnesses just chosen, and the inverse of $g(0,c)\Delta(0,c)$.
This is a nonzero finitely generated $\ZZ$-algebra of characteristic
zero.  The same presentation defines an etale $W[t]$-algebra
\[
 E_W=\bigl(W[t,z]/(G_1,\ldots,G_m)\bigr)_{g\Delta}.
\]
Evaluation at $(0,c)$ gives a $W$-valued point.  All solution
identities hold in $E_W$, by the chosen witnesses.

Choose a maximal ideal of $W$ and let $k_0$ be its residue field.
This field is finite.  Indeed, it is a field finitely generated as a
$\ZZ$-algebra.  If it had characteristic zero, Zariski's lemma over
$\mathbb Q$ would make it a number field.  Finitely many ring
generators of a number field are contained in its ring of integers
with one nonzero integer inverted, which cannot contain the inverse
of every rational prime.  This is impossible for a field of
characteristic zero.  In positive characteristic, Zariski's lemma
makes $k_0$ a finite extension of its finite prime field.

Let $K$ be an algebraic closure of $k_0$.  The distinguished point
specializes to a $K$-rational point of $E_W\otimes_W K$ over $t=0$.
The inverted Jacobian and denominator remain nonzero there.
The etale lifting property, or multivariable Hensel's lemma, lifts
this point to a $K[t]$-algebra map
\[
 E_W\otimes_W K\longrightarrow K[[t]].
\]
Explicitly, the point lifts uniquely over each successive quotient
$K[[t]]/(t)^a$, since the Jacobian is invertible, and the compatible
lifts give the displayed map by completeness.  Applying it to the
solution entries preserves all the polynomial identities and proves
the assertion.
\end{proof}

\begin{theorem}\label{thm:charzero}
For every flat local homomorphism of nonzero Noetherian local rings
$(R,\mathfrak m)\to(S,\mathfrak n)$ whose residue fields have
characteristic zero, one has $e(R)\le e(S)$.
\end{theorem}

\begin{proof}
By Lemma~\ref{lem:reduction}, it suffices to consider complete rings
of the same dimension with algebraically closed target residue field.
Lemma~\ref{lem:charzero-finite-free} then permits us to assume that
the map is finite free and that both rings have the same
algebraically closed coefficient field $L$.

If the common dimension is zero, equality of residue fields and
freeness of rank $n\ge1$ give
$e(S)=\length_S S=n\length_R R\ge e(R)$.
Suppose instead that $d>0$ and that $e(R)>e(S)$.
Lemma~\ref{lem:charzero-certificate} produces a finite polynomial
system with a solution over $L[[t_1,\ldots,t_d]]$.
Lemma~\ref{lem:charzero-specialization} gives a solution over
$K[[t_1,\ldots,t_d]]$ for an algebraically closed field $K$ of
positive characteristic.  The resulting finite free local map
$\widetilde R\to\widetilde S$ satisfies
\[
 e(\widetilde R)\ge e(R)>e(S)\ge e(\widetilde S),
\]
contradicting Theorem~\ref{thm:positive-characteristic}.
\end{proof}
